\subsection{Fourier conventions and the local reduction}
\label{red:section}

We use the Fourier convention
\[
 \widehat f(\xi,\eta)
   =\int_{\R^2}f(x,y)e^{-i(x\xi+y\eta)}\dd x\dd y,
 \qquad
 f(x,y)=\frac1{(2\pi)^2}\int_{\R^2}
       \widehat f(\xi,\eta)e^{i(x\xi+y\eta)}\dd\xi\dd\eta.
\]
The same convention is used in one dimension. Pairings are linear in
the first entry. A constant depending on a fixed finite list of smooth
cutoffs or on their prescribed seminorms is an absolute constant.
No such list will depend on a frequency scale, an inner truncation,
the number of packets, or the dimension of the bin Hilbert space.

\subsection{A maximal theorem and elementary analytic tools}

We first record precisely the specialized maximal theorem used later.
If an orthogonal rectangle \(R\) has long side \(L(R)\), short side
\(W(R)\), and long-side direction \(e_R\in S^1\), let
\(\operatorname{EX}(R)\) be the arc centered at \(e_R\) of total
length \(W(R)/L(R)\).

\begin{theorem}[Lacey--Li]\label{red:kakeya}
Let \(v:\R^2\to S^1\) be Lipschitz and \(0<\delta<1\). Restrict to
rectangles satisfying
\[
 L(R)\le \frac1{100\Lip(v)},\qquad
 \abs{R\cap v^{-1}(\operatorname{EX}(R))}\ge\delta\abs R,
\]
where the first restriction is void when \(\Lip(v)=0\). Then
\[
 \mathcal M_{v,\delta}h(z)
   =\sup_R\frac{\one_R(z)}{\abs R}\int_R\abs h
 \quad\text{satisfies}\quad
 \norm{\mathcal M_{v,\delta}h}_{L^{2,\infty}}
       \le C\delta^{-1/2}\norm h_2.
\]
Here the supremum is tested at every point of \(R\), and the constant
is independent of \(v\) and \(\delta\).
\end{theorem}

This is \cite[Theorem~1.4]{LLKakeya}. Requiring the field to lie in a
smaller fixed subarc only restricts the collection of admissible
rectangles. In the bounded slope range, the map
\(t\mapsto(1,t)/\sqrt{1+t^2}\) changes slope distances into arc
distances by comparable absolute factors. These observations explain
the slope version of \Cref{red:kakeya} used below. When an intermediate
box is a parallelogram, the application of the theorem will explicitly
enclose it by a comparable orthogonal rectangle of admissible length.

We also use the following familiar one-dimensional estimates, with
their indicated uniformities. The Hardy--Littlewood maximal operator
is bounded on \(L^p\), \(1<p<\infty\), and
\begin{equation}\label{red:horizontal-hilbert}
 \sup_{0<r<R<\infty}
 \left\|\int_{r<|t|<R}h(\,\cdot-t\,)\frac{\dd t}{t}\right\|_p
   \le C_p\norm h_p.
\end{equation}
The supremum in \eqref{red:horizontal-hilbert} is outside the norm;
the usual maximal truncation theorem gives this assertion as well.
See \cite[Theorems~1.1 and~7.6 and the proof of (7.7)]{Hajlasz2012}.
Fubini gives the corresponding horizontal estimates on \(\R^2\).
The \(L^2\) square-function identities used here are consequences of
Plancherel and finite overlap of smooth frequency annuli. The standard
smooth Littlewood--Paley and random-sign multiplier formulations are
also recorded in \cite[Proposition~5.3, Corollary~5.4, and
Remark~4.5]{Tao2006Notes4}. We give the particular Carleson embedding
and commutator arguments needed for our reduction below. The
John--Nirenberg and martingale variation estimates are stated in the
sections where their local forms are used.

\subsection{The local slope operator and its endpoints}

Fix an absolute bound \(C_{\rm Lip}\) large enough for the slope
fields obtained in the following localization; for example, a bound
of ten suffices after taking the squares small enough. Choose an
absolute horizontal length \(b_0>0\) with
\(b_0 C_{\rm Lip}<1/10\), and subsequently choose
\(0<a_*<b_0/100\). These constants may be decreased to accommodate
the fixed geometric enlargements in later applications of
\Cref{red:kakeya}. All such enlargements are chosen absolutely.
For the remainder of the local argument write \(a=a_*\), and put
\begin{equation}\label{red:local-operator}
 H_{u,a}^{\epsilon}h(x,y)
   =\int_{\epsilon<|t|<a}
      h(x-t,y-tu(x,y))\frac{\dd t}{t},
 \qquad 0<\epsilon<a,
\end{equation}
where \(\abs u\le1\) and \(\Lip(u)\le C_{\rm Lip}\).
The local inputs and tests will be supported in
\(I\times\R\), where \(I\) is any interval of length at most
\(b_0\).

\begin{lemma}\label{red:fixed-time}
For \(\abs t\,C_{\rm Lip}<1\), define
\[
 A_t h(x,y)=h(x-t,y-tu(x,y)).
\]
For \(1\le p<\infty\),
\begin{equation}\label{red:translate-bound}
 \norm{A_t h}_p
   \le (1-\abs t\,C_{\rm Lip})^{-1/p}\norm h_p,
\end{equation}
and \(A_t\) is a contraction on \(L^\infty\).
Consequently, replacing either symmetric endpoint at scale \(r\)
by a measurable endpoint between \(c r\) and \(C r\), with
\(c,C>0\) fixed and \(\max(1,C)r C_{\rm Lip}<1/2\), costs a bounded
\(L^p\) operator, independently of \(r\).
\end{lemma}

\begin{proof}
For each fixed \(x\), the map
\(\varphi_{x,t}(y)=y-tu(x,y)\) is increasing, onto, and satisfies
\[
 (1-\abs t\,C_{\rm Lip})\abs{y-y'}
 \le\abs{\varphi_{x,t}(y)-\varphi_{x,t}(y')}
 \le(1+\abs t\,C_{\rm Lip})\abs{y-y'}.
\]
The change-of-variables formula for a one-dimensional bi-Lipschitz
map, followed by translation in \(x\), proves
\eqref{red:translate-bound}. A difference of endpoint truncations
is pointwise dominated by
\[
 \int_{c' r<|t|<C' r}\abs{A_t h}\frac{\dd t}{|t|}
\]
with \(c'=\min(1,c)\) and \(C'=\max(1,C)\). Minkowski and
\eqref{red:translate-bound} bound this by
\(C_p\log(C'/c')\norm h_p\). The same proof applies to inner
and outer endpoints, and permits endpoints depending on \((x,y)\).
\end{proof}

We explain the localization of the original unit-field operator.
Partition the plane into squares \(Q\) of a sufficiently small fixed
side length. For output in \(Q\), only a fixed enlargement \(Q^+\)
of that square contributes to a transform of length \(a_*\). Rotate
coordinates so that the vector at the center of \(Q\) points along
the positive first coordinate axis. Since \(\Lip(v)\le1\), on the
rotated output square we have
\[
 v=(v_1,v_2),\qquad v_1\ge c>0,
 \qquad u=v_2/v_1,\qquad \abs u\le1,
 \qquad\Lip(u)\le C_{\rm Lip}.
\]
The last inequality follows by subtracting the two ratios and using
the lower bound for \(v_1\). Extend \(u\) to the plane with the
same Lipschitz bound, using
\(\inf_{z\in Q}(u(z)+C_{\rm Lip}|\,\cdot-z|)\), and then clip
the extension to \([-1,1]\). The infimum extension agrees with
\(u\) on \(Q\), and clipping preserves both agreement and the
Lipschitz bound.

At an output point, the substitution \(s=t v_1(x,y)\) changes
\(\dd t/t\) to \(\dd s/s\) and changes the two endpoints to
\(\epsilon v_1(x,y)\) and \(a_*v_1(x,y)\). By
\Cref{red:fixed-time}, each endpoint can be replaced by
\(\epsilon\) or \(a_*\), respectively, at a uniform \(L^2\)
cost. The rotated sets \(Q\) and \(Q^+\) fit in a common
horizontal interval of length at most \(b_0\), after choosing the
square side and \(a_*\) in that order. Thus a uniform bound for
\(\one_{I\times\R}H_{u,a}^{\epsilon}\one_{I\times\R}\)
implies the required bound on this output square. The output squares
are disjoint, the sets \(Q^+\) have bounded overlap, and rotations
preserve area. Squaring the local estimates and summing over \(Q\)
therefore proves the global estimate. Notice that both endpoint
errors are uniform in the inner radius \(\epsilon\).

\subsection{The band theorem and summation of amplitude levels}

A vertical band cutoff at length \(w>0\) is a multiplier
\(P_w=p(wD_y)\), where \(p\) is smooth and supported in a fixed
compact annulus
\(\{r:c_\eta\le|r|\le C_\eta\}\), with
\(0<c_\eta<C_\eta<\infty\) absolute. We allow a fixed bounded
family of such cutoffs, including slightly enlarged cutoffs equal to
one on the supports of others. All constants below are uniform when
the required finitely many cutoff seminorms are bounded. The dyadic
inverse frequency scale \(w\) belongs to \(2^{\Z}\).

\begin{theorem}[Restricted estimate on a vertical band]
\label{thm:band-restricted}
For the fixed absolute constants in \eqref{red:local-operator}, there
is an absolute \(C\) such that, for every field \(u\) with
\(\abs u\le1\) and \(\Lip(u)\le C_{\rm Lip}\), every
\(0<\epsilon<a\), every dyadic \(w\), and every vertical band
cutoff \(P_w\),
\begin{equation}\label{eq:src-1}
 \abs{\langle H_{u,a}^{\epsilon}P_w f,g\rangle}
 \le C\bigl(1+\abs{\log(|F|/|E|)}\bigr)^{-3}
       \sqrt{|F||E|}
\end{equation}
whenever \(F,E\subset I\times\R\) have positive finite measure,
\(\abs f\le\one_F\), \(\abs g\le\one_E\), and
\(\abs I\le b_0\).
\end{theorem}

The proof of \Cref{thm:band-restricted} will be completed after the
packet estimates and their iteration have been established. The next
results show, without using that theorem as a premise elsewhere in
its proof, that it implies \Cref{thm:main}. The last part of this
section then gives explicit sufficient packet estimates for
\Cref{thm:band-restricted}.

\begin{lemma}\label{red:restricted-strong}
Suppose a bilinear form \(\mathcal B\) satisfies
\[
 \abs{\mathcal B(f,g)}
 \le C_0(1+|\log(|F|/|E|)|)^{-3}\sqrt{|F||E|}
 \quad (|f|\le\one_F,\ |g|\le\one_E)
\]
on finite-measure sets in a fixed measurable domain. Then
\[
 \abs{\mathcal B(f,g)}\le C C_0\norm f_2\norm g_2
\]
for finite amplitude decompositions, and hence for all \(L^2\)
inputs whenever the form is obtained by an \(L^2\)-continuous
truncation or by completion from those inputs.
\end{lemma}

\begin{proof}
Use disjoint amplitude sets
\(F_i=\{2^{i-1}<|f|\le2^i\}\) and write
\(f=\sum_i2^i f_i\), where \(|f_i|\le\one_{F_i}\).
For the moment all sums are finite. Put
\(a_i=2^i\sqrt{|F_i|}\), so
\(\sum_i a_i^2\le4\norm f_2^2\).
Group the indices into measure classes
\[
 \mathcal I_k=\{i:2^k\le|F_i|<2^{k+1}\},\qquad
 A_k=\left(\sum_{i\in\mathcal I_k}a_i^2\right)^{1/2}.
\]
If \(i_0\) is the largest index in a nonempty class, then
\[
 \sum_{i\in\mathcal I_k}a_i
 \le 2^{(k+1)/2}\sum_{i\le i_0}2^i
 \le 2\sqrt2\,a_{i_0}\le2\sqrt2\,A_k.
\]
Make the analogous decomposition \(g=\sum_j2^j g_j\), with
measure classes \(\mathcal J_l\) and square sums \(B_l\).
For \(i\in\mathcal I_k\), \(j\in\mathcal J_l\),
\[
 (1+|\log(|F_i|/|E_j|)|)^{-3}
       \le C(1+|k-l|)^{-3}.
\]
The assumed bound, followed by the within-class estimates, gives
\[
 \abs{\mathcal B(f,g)}
   \le C C_0\sum_{k,l}(1+|k-l|)^{-3}A_kB_l
   \le C C_0\norm A_{\ell^2}\norm B_{\ell^2}.
\]
The final step is Young's inequality for the summable sequence
\((1+|k|)^{-3}\). This proves the assertion. Truncating the amplitude
decompositions and using the asserted continuity or completion gives
the last statement, with no dependence on the number of classes.
\end{proof}

\subsection{A Lipschitz commutator}

The estimate below belongs to the Lipschitz commutator theory initiated
by Calder\'on and developed by Coifman and Meyer
\cite{Calderon1965,CoifmanMeyer1978}. We give a dyadic proof with the
precise multiplier seminorms needed for Littlewood--Paley reassembly
and for the later narrow-bin reduction. The argument separates low
output frequencies from the annular interactions, because the constant
must remain uniform as the number of frequency scales increases.
Let \(\chi\) be a real even smooth cutoff equal to one on
\([-1,1]\) and supported in \([-2,2]\), and define
\[
 S_k=\chi(2^{-k}D),\qquad \Delta_k=S_k-S_{k-1}.
\]
Thus \(\Delta_k\) is supported in
\(\{2^{k-1}\le|\zeta|\le2^{k+1}\}\), and Plancherel gives
\begin{equation}\label{red:lp-square}
 \norm h_2^2\asymp\sum_k\norm{\Delta_kh}_2^2.
\end{equation}

\begin{lemma}[The required Carleson embedding]\label{red:carleson}
For every fixed integer \(c\), every \(q\in L^\infty(\R)\), and
every \(h\in L^2(\R)\),
\begin{equation}\label{red:carleson-estimate}
 \sum_k\norm{(\Delta_kq)S_{k+c}h}_2^2
      \le C_c\norm q_\infty^2\norm h_2^2.
\end{equation}
\end{lemma}

\begin{proof}
For a dyadic interval \(J\) of length \(2^{-j}\), split
\(q=q\one_{3J}+q\one_{(3J)^c}\). By
\eqref{red:lp-square}, the first part satisfies
\[
 \sum_{k\ge j}\int_J
       |\Delta_k(q\one_{3J})|^2
       \le C\norm q_\infty^2|J|.
\]
The kernels of \(\Delta_k\) have Schwartz decay at scale
\(2^{-k}\). For \(x\in J\) and \(k\ge j\), this gives
\[
 |\Delta_k(q\one_{(3J)^c})(x)|
     \le C\norm q_\infty 2^{j-k}.
\]
After summing the squares, we obtain
\begin{equation}\label{red:carleson-measure}
 \sum_{k\ge j}\int_J|\Delta_kq|^2
          \le C\norm q_\infty^2|J|.
\end{equation}
For each dyadic interval \(I\) of length \(2^{-k}\), set
\(\alpha_I=\int_I|\Delta_kq|^2\). Then
\eqref{red:carleson-measure} says
\(\sum_{I\subset J}\alpha_I\le C\norm q_\infty^2|J|\).

Write \(M\) for the Hardy--Littlewood maximal operator. Kernel
decay, summed on dyadic annuli about a point of \(I\), gives
\[
 \sup_{x\in I}|S_{k+c}h(x)|
       \le C_c\operatorname*{ess\,inf}_{y\in I}Mh(y).
\]
The shift \(c\) is fixed, so moving the kernel center across \(I\)
changes the bound by only a fixed factor. If \(F\ge0\) belongs to
\(L^2\), the dyadic Carleson condition implies
\[
 \sum_I\alpha_I
     \bigl(\operatorname*{ess\,inf}_I F\bigr)^2
       \le C\norm q_\infty^2\norm F_2^2.
\]
Indeed, for each \(t>0\), sum over the maximal dyadic intervals
on which \(\operatorname*{ess\,inf}_I F>t\), apply the Carleson
condition on each such interval, and integrate against \(2t\dd t\).
Their union lies, up to null sets, in \(\{F>t\}\).
Apply this assertion with \(F=Mh\), and use the \(L^2\) maximal
theorem. This proves \eqref{red:carleson-estimate}. This is also the
usual continuous Littlewood--Paley Carleson embedding; compare the
proof of \cite[Proposition~4.1]{Tao2007Notes6}.
\end{proof}

\begin{lemma}\label{red:commutator}
Let \(b\in W^{1,\infty}(\R)\), and let \(T=\mathfrak m(D)\),
where \(\mathfrak m\) is smooth away from zero and
\[
 A(\mathfrak m)
   =\max_{0\le r\le4}\sup_{\zeta\ne0}
        |\zeta|^r|\mathfrak m^{(r)}(\zeta)|<\infty.
\]
Then
\begin{equation}\label{red:commutator-estimate}
 \norm{[b,T]\partial h}_2
    \le C A(\mathfrak m)\norm{b'}_\infty\norm h_2.
\end{equation}
The commutator is initially defined on the Sobolev space
\(W^{1,2}(\R)\); the estimate gives its bounded extension to
\(L^2\). Its constant is independent
of \(\norm b_\infty\) and of all frequency truncations.
\end{lemma}

\begin{proof}
Put \(L_b=\norm{b'}_\infty\) and \(b_k=\Delta_kb\).
Cancellation and the first absolute moment of the smooth kernels give
\begin{equation}\label{red:b-band}
 \norm{b_k}_\infty\le C2^{-k}L_b,
 \qquad \norm{\partial S_kb}_\infty\le C L_b,
 \qquad b_k'=\Delta_kb'.
\end{equation}
Also \(\norm{S_Nb-b}_\infty\le C2^{-N}L_b\).
Consequently, for each fixed \(j\), the exact identity
\begin{equation}\label{red:b-reconstruction}
 b=S_{j-4}b+\sum_{k\ge j-3}b_k
\end{equation}
has a uniformly convergent high-frequency sum. This identity retains
the constant and low-frequency part of \(b\); no homogeneous
reconstruction of a bounded function at frequency zero is needed.

For any fixed smooth annular cutoff \(\rho\), let \(K_j\) be the
kernel of \(\mathfrak m(D)\rho(2^{-j}D)\). Rescaling and three
integrations by parts on its fixed compact frequency support show
\begin{equation}\label{red:kernel-first-moment}
 \int_\R |z|\,|K_j(z)|\dd z
       \le C A(\mathfrak m)2^{-j}.
\end{equation}
In fact the rescaled kernel is bounded by
\(C A(\mathfrak m)(1+|z|)^{-3}\), which proves the assertion.

First suppose that \(\widehat h\) is smooth and compactly supported
away from zero. Put \(h_j=\Delta_jh\); only finitely many \(h_j\)
are nonzero. Apply \eqref{red:b-reconstruction} to each of them.
There are three interactions.

For the low frequencies of \(b\), put \(B_j=S_{j-4}b\).
The supports of both \(h_j\) and \(B_jh_j\) lie in a fixed
annulus of radius \(2^j\). A slightly enlarged annular multiplier
\(T_j\) therefore satisfies
\([B_j,T]\partial h_j=[B_j,T_j]\partial h_j\).
The kernel formula and \eqref{red:kernel-first-moment} give
\[
 \norm{[B_j,T_j]\partial h_j}_2
 \le \Lip(B_j)\int |z||K_j(z)|\dd z\,
             \norm{\partial h_j}_2
 \le C A(\mathfrak m)L_b\norm{h_j}_2.
\]
The outputs remain in fixed annuli of radius \(2^j\). Finite
frequency overlap and \eqref{red:lp-square} bound their sum by
\(C A(\mathfrak m)L_b\norm h_2\).

For \(k\ge j+4\), the output of
\([b_k,T]\partial h_j\) lies in a fixed annulus of radius \(2^k\).
Plancherel, Bernstein's inequality, and \eqref{red:b-band} yield
\[
 \norm{[b_k,T]\partial h_j}_2
       \le C A(\mathfrak m)L_b 2^{j-k}\norm{h_j}_2.
\]
Group by \(k\), use frequency overlap, and apply Young's inequality
on \(\ell^2(\Z)\) to the summable sequence
\(2^{-n}\one_{n\ge4}\). This bounds the high--low sum by the
same quantity and proves its convergence in \(L^2\).

It remains to consider \(|k-j|\le3\). Fix one such difference and
write \(H_k=h_{k+r}\), \(|r|\le3\). The exact identity
\begin{equation}\label{red:commutator-high-high}
 [b_k,T]\partial H_k
  =\partial([b_k,T]H_k)-b_k'TH_k+T(b_k'H_k)
\end{equation}
treats the possible low output frequencies. Indeed,
\(\partial([b_k,T]H_k)\) has spectrum in a ball of radius
\(C2^k\). Its \(\Delta_i\) projection vanishes for
\(i>k+c_0\), with \(c_0\) fixed, and otherwise
\[
 \norm{\Delta_i\partial([b_k,T]H_k)}_2
    \le C A(\mathfrak m)L_b2^{i-k}\norm{H_k}_2.
\]
Another application of \eqref{red:lp-square} and Young's inequality
to \(2^{-n}\one_{n\ge-c_0}\) bounds the sum of these derivative
terms. Thus arbitrarily low output frequencies cause no logarithmic
loss.

Choose a fixed \(c_1\) so that \(S_{k+c_1}\) equals one on the
spectrum of every product \(b_k'H_k\). For \(g\in L^2\),
\[
 \begin{split}
 \left|\left\langle\sum_k b_k'H_k,g\right\rangle\right|
 &\le
 \left(\sum_k\norm{H_k}_2^2\right)^{1/2}
 \left(\sum_k\norm{(\Delta_kb')S_{k+c_1}g}_2^2\right)^{1/2}\\
 &\le C L_b
       \left(\sum_k\norm{H_k}_2^2\right)^{1/2}\norm g_2,
 \end{split}
\]
by \Cref{red:carleson}. Since \(\norm T_{2\to2}\le A(\mathfrak m)\),
this bounds \(T\sum_k b_k'H_k\). Replacing \(H_k\) by \(TH_k\)
in the same calculation bounds \(\sum_k b_k'TH_k\): the frequency
supports are unchanged and
\(\sum_k\norm{TH_k}_2^2\le A(\mathfrak m)^2
\sum_k\norm{H_k}_2^2\). Equation~\eqref{red:commutator-high-high},
summed over its seven possible differences, completes the proof on
the chosen dense class.

For \(h\in W^{1,2}\), approximate by
\((S_N-S_{-N})h\), followed if necessary by smooth compact Fourier
approximations away from zero. The approximation converges in
\(W^{1,2}\). Since \(b\) is bounded and \(T\) is \(L^2\)-bounded,
the defining commutators converge in \(L^2\). The estimate passes
to the limit, and then defines a bounded operator on all of \(L^2\).
\end{proof}

For a vertical multiplier \(T\), set
\[
 C_t h(x,y)=h(x,y-tu(x,y)).
\]

\begin{corollary}\label{red:composition-commutator}
For \(\abs t\le b_0\) and \(T\) as in \Cref{red:commutator},
\begin{equation}\label{red:composition-bound}
 \norm{TC_t h-C_tTh}_2
      \le C A(\mathfrak m)\abs t\,\norm h_2.
\end{equation}
The constant is uniform over \(\abs u\le1\),
\(\Lip(u)\le C_{\rm Lip}\).
\end{corollary}

\begin{proof}
Work first with a smooth field and on one fixed \(x\)-fiber.
Let \(\varphi_s(y)=y-su(x,y)\). Direct differentiation gives
\[
 C_s^{-1}C_s'=b_s\partial_y,
 \qquad b_s=-u(x,\cdot)\circ\varphi_s^{-1},
\]
and hence, since \(T\) commutes with \(\partial_y\),
\begin{equation}\label{red:conjugation-derivative}
 \frac{\dd}{\dd s}(C_sTC_s^{-1})
       =C_s[b_s,T]\partial_y C_s^{-1}.
\end{equation}
The bi-Lipschitz estimates in \Cref{red:fixed-time} imply
\[
 \norm{C_s}_{2\to2}+\norm{C_s^{-1}}_{2\to2}\le C,
 \qquad
 \Lip(b_s)\le
    \frac{C_{\rm Lip}}{1-|s|C_{\rm Lip}}\le C.
\]
Apply \Cref{red:commutator} inside
\eqref{red:conjugation-derivative}, integrate from \(0\) to \(t\),
and multiply by \(C_t\). The estimate is uniform in the fiber, so
integration in \(x\) proves \eqref{red:composition-bound}.

A bounded Lipschitz field is approximated by convolutions with smooth
approximate identities, preserving the bounds on its size and
Lipschitz seminorm. For each fixed small \(s\), the corresponding
composition operators converge strongly on \(L^2\): this follows
first for continuous compactly supported inputs by uniform
convergence of the maps, and then by density and their uniform
operator bounds. Passing to the limit proves the assertion for the
original field.
\end{proof}

\subsection{Square-function reassembly}

The preceding composition estimate pays a factor of \(|t|\), which
cancels the singular measure \(\dd t/|t|\) when commutators are
integrated. This is the mechanism that permits the band estimates to
be assembled. Di Plinio, Guo, Thiele and Zorin-Kranich established
this band-to-full reduction under Lipschitz regularity
\cite[Corollaries~1.3--1.4 of the author version]{DiPlinioGuoThieleZorinKranich2018}.
We retain the argument here to track the fixed outer endpoint,
arbitrary inner endpoint and horizontal localization used throughout
this paper.


Choose a real smooth annular function \(p\) so that
\(\sum_{w\in2^{\Z}}p(w\eta)^2=1\) for \(\eta\ne0\).
Such a function is obtained by dividing an ordinary smooth dyadic
partition by the square root of its sum of squares. Set
\(P_w=p(wD_y)\). For every finite set \(\mathcal W\) of scales and
every choice of signs \(\varepsilon_w\in\{-1,1\}\), the multiplier
\(T=\sum_{w\in\mathcal W}\varepsilon_wP_w\) satisfies
\(A(\mathfrak m_T)\le C\), by finite overlap of the annuli.
Writing \(S_t^x h(x,y)=h(x-t,y)\), we have
\[
 H_{u,a}^{\epsilon}
      =\int_{\epsilon<|t|<a} C_t S_t^x\frac{\dd t}{t}.
\]
The vertical multiplier \(T\) commutes with \(S_t^x\). Thus
\Cref{red:composition-commutator} implies
\begin{equation}\label{red:integrated-commutator}
 \norm{[T,H_{u,a}^{\epsilon}]}_{2\to2}
       \le C\int_{\epsilon<|t|<a}|t|\frac{\dd t}{|t|}
       \le C a.
\end{equation}
All these operators are already bounded on \(L^2\) for fixed
\(\epsilon>0\), with the elementary bound
\(C\log(a/\epsilon)\); no conclusion of
\Cref{thm:band-restricted} is needed to form their square functions.

Randomize the signs independently and use their orthogonality.
For \(f=\one_{I\times\R}f\), Equation~\eqref{red:integrated-commutator}
and the triangle inequality in the random-sign \(L^2\) space give
\begin{equation}\label{red:reassembly}
 \left(\sum_{w\in\mathcal W}
       \norm{\one_{I\times\R}P_w H_{u,a}^{\epsilon}f}_2^2
       \right)^{1/2}
 \le C\norm f_2+
 \left(\sum_{w\in\mathcal W}
       \norm{\one_{I\times\R}H_{u,a}^{\epsilon}P_w f}_2^2
       \right)^{1/2}.
\end{equation}
The horizontal restriction commutes with every vertical multiplier.

\begin{proposition}\label{red:band-implies-main}
The uniform restricted estimate in \Cref{thm:band-restricted} implies
\Cref{thm:main}.
\end{proposition}

\begin{proof}
Apply \Cref{red:restricted-strong} to the localized band form.
It gives a uniform local \(L^2\) bound for
\(H_{u,a}^{\epsilon}\widetilde P_w\), where
\(\widetilde P_w\) is an enlarged cutoff equal to one on the
spectrum of \(P_w\). Since \(P_w f\) retains the same horizontal
support,
\[
 \norm{\one_{I\times\R}H_{u,a}^{\epsilon}P_w f}_2
 =\norm{\one_{I\times\R}H_{u,a}^{\epsilon}
                      \widetilde P_w(P_w f)}_2
 \le C\norm{P_w f}_2.
\]
Insert this in \eqref{red:reassembly} and let \(\mathcal W\)
increase to all dyadic scales. Plancherel and the square partition
give
\(\norm{\one_{I\times\R}H_{u,a}^{\epsilon}f}_2
\le C\norm f_2\), uniformly in \(\epsilon\). The localization
already proved above gives the asserted global bound for the unit
field. Finally, for a Schwartz input the symmetric principal value
exists pointwise; Fatou gives its \(L^2\) bound, and Chebyshev gives
the weak estimate in \Cref{thm:main}.
\end{proof}

\subsection{Removing long and very short times}

We now reduce \Cref{thm:band-restricted} to packet estimates. In this
part only, abbreviate \(H_{u,a}^{\epsilon}\) to \(H\), and put
\[
 L=\log\sqrt{|E|/|F|}.
\]
Fix a sufficiently large absolute threshold \(L_0\), which can be
increased after the fixed constants in the argument are chosen.
When \(L\ge L_0\), choose dyadic numbers satisfying
\begin{equation}\label{red:small-time-parameters}
 \tfrac12 e^{-L^{0.97}}<\tau\le e^{-L^{0.97}},\qquad
 e^{L^{0.9}}\le m<2e^{L^{0.9}},\qquad W=mw.
\end{equation}
Here \(\tau\) is the short-time cutoff. The parameter \(m\)
sets the number of vertical frequency bins: their width will be
comparable to \(W^{-1}\), and \(W=mw\) is the corresponding
spatial width. For every fixed \(C\), the product \(m^C\tau\)
tends to zero as \(L\to\infty\); this pays for the binning errors.
We take \(L_0\) large enough that \(\tau<a\).
For all other values of \(L\), let \(m=1\), \(W=w\), and let
\(\tau\in(0,a)\) be a fixed small dyadic constant.

\begin{lemma}\label{red:time-reduction}
For any fixed \(1<p<\infty\), the contribution of
\(\max(\epsilon,\tau)<|t|<a\) to \(HP_w\) has \(L^p\) norm
at most \(C_p(1+L)\) if \(L\ge L_0\), and at most a fixed
\(C_p\) otherwise. On the remaining times below \(w\), replacing
\(u\) by zero costs a uniformly bounded \(L^p\) operator. The
resulting horizontal symmetric truncation is uniformly bounded on
\(L^p\).
\end{lemma}

\begin{proof}
The first assertion follows from \Cref{red:fixed-time}, the uniform
\(L^1\) norm of the convolution kernel of \(P_w\), and
\[
 \int_{\tau<|t|<a}\frac{\dd t}{|t|}
       =2\log(a/\tau)\le C(1+L^{0.97}).
\]
For the second, the fundamental theorem of calculus gives
\[
 \begin{split}
 &(P_wh)(x-t,y-tu(x,y))-(P_wh)(x-t,y)\\
 &\qquad=-t u(x,y)\int_0^1
    (\partial_yP_wh)(x-t,y-\theta t u(x,y))\dd\theta.
 \end{split}
\]
The intermediate maps satisfy \Cref{red:fixed-time}, and the kernel
of \(\partial_yP_w\) has \(L^1\) norm \(O(w^{-1})\).
Consequently,
\begin{equation}\label{red:short-time-error}
 \norm{(A_t-S_t^x)P_wh}_p
        \le C_p |t|w^{-1}\norm h_p.
\end{equation}
Integrating over any subinterval of \(0<|t|<w\) against
\(\dd t/|t|\) gives a constant. Equation~\eqref{red:horizontal-hilbert}
handles the horizontal part, uniformly in its inner and outer
endpoints.
\end{proof}

Every error with \(L^p\) norm \(B_p\) contributes at most
\begin{equation}\label{red:restricted-error}
 B_p|F|^{1/p}|E|^{1-1/p}
   =B_p\sqrt{|F||E|}
           \left(\frac{|F|}{|E|}\right)^{1/p-1/2}
\end{equation}
to the restricted form. In the regime \(L\ge L_0\), choose one
fixed \(p<2\), for example \(p=3/2\); the normalized error is
then at most \(C(1+L)e^{-L/3}\). In the opposite regime
\(|F|\ge|E|\), the fixed choice \(p=4\) gives the power
\((|E|/|F|)^{1/4}\) for each uniformly bounded error. On a bounded
range of \(L\), fixed operator bounds suffice. These choices are
made once and do not approach the endpoint \(p=2\).

\subsection{From a spatial truncation to a finite scale sum}

Choose an even smooth cutoff \(\chi_0\), equal to one near zero
and compactly supported, and set
\[
 \vartheta(r)=\chi_0(r)-\chi_0(2r),\qquad
 \psi(r)=-i\pi\operatorname{sgn}(r)\vartheta(r).
\]
The function \(\psi\) is smooth and supported in a fixed compact
annulus. Its inverse Fourier transform \(k\) is Schwartz, and
\(k_\ell(t)=\ell^{-1}k(t/\ell)\) has multiplier
\(\psi(\ell\zeta)\). Since
\(\sum_{\ell\in2^{\Z}}\vartheta(\ell\zeta)=1\) for
\(\zeta\ne0\),
\[
 \sum_{\ell\in2^{\Z}}k_\ell(t)=\frac1t\qquad(t\ne0).
\]
The identity follows by Fourier inversion in tempered distributions;
away from zero the kernel sum is locally uniformly convergent by
Schwartz decay for small \(\ell\) and the bound
\(C/\ell\) for large \(\ell\).

\begin{lemma}\label{red:scale-truncation}
For dyadic \(0<r\le R\),
\begin{equation}\label{red:scale-kernel-error}
 \int_\R\left|
    \frac{\one_{\{r<|t|<R\}}}{t}
           -\sum_{r\le\ell\le R}k_\ell(t)
       \right|\dd t\le C.
\end{equation}
Also,
\begin{equation}\label{red:scale-first-moment}
 \int_\R |t|\left|\sum_{r\le\ell\le R}k_\ell(t)\right|
        \dd t\le C R.
\end{equation}
Both constants are independent of the number of scales. Changing
either spatial endpoint by a fixed factor preserves the first bound.
\end{lemma}

\begin{proof}
For \(r<|t|<R\), the omitted scales below \(r\) have sum bounded
by \(C r/|t|^2\), using
\(\abs{k_\ell(t)}\le C\ell/|t|^2\) and a geometric sum.
The omitted scales above \(R\) have sum at most \(C/R\).
The integrals of these bounds over the middle interval are bounded.
For \(|t|<r\), the retained sum is at most
\(C\sum_{\ell\ge r}\ell^{-1}\le C/r\); for \(|t|>R\),
it is at most \(C R/|t|^2\). These are integrable on their
respective regions and prove \eqref{red:scale-kernel-error}.
For the moment bound, use
\[
 \int |t||k_\ell(t)|\dd t
          =\ell\int |s||k(s)|\dd s
\]
and \(\sum_{\ell\le R}\ell\le2R\). Finally, changing an
endpoint by a fixed factor changes the spatial kernel only on an
annulus having bounded \(\dd t/|t|\) measure.
\end{proof}

After \Cref{red:time-reduction}, the only possible remaining times
are between \(r=\max(w,\epsilon)\) and \(\tau\). Round the
lower endpoint to a comparable dyadic number; the resulting annular
error has a uniform \(L^p\) bound. If this lower endpoint exceeds
\(\tau\), there is no remaining model. Otherwise
\Cref{red:scale-truncation} reduces the multiplier, with \(u\)
frozen at the output point, to
\begin{equation}\label{eq:src-2}
 \mathfrak a_{u,\mathcal L}(x,y;\xi,\eta)
    =\sum_{\ell\in\mathcal L}
           \psi\bigl(\ell(\xi+\eta u(x,y))\bigr),
 \qquad
 \mathcal L\subset\{\ell\in2^{\Z}:w\le\ell\le\tau\}.
\end{equation}
Here \(\mathcal L\) is the consecutive set of retained scales;
when \(\epsilon\le w\) it is the whole displayed range. Denote
the associated operator by \(\mathcal A_{u,\mathcal L}\), or
simply \(\mathcal A\). Its spatial representation is
\[
 \mathcal A h(x,y)
   =\int_\R K_{\mathcal L}(t)
                  h(x-t,y-tu(x,y))\dd t,
 \qquad K_{\mathcal L}=\sum_{\ell\in\mathcal L}k_\ell.
\]

The kernels in this representation are not compactly supported in
\(t\). This causes no large-time change-of-variables problem: if
both the input and the test are supported in \(I\times\R\), a
nonzero contribution has \(t\in I-I\), hence \(|t|\le b_0\).
Every vertical cutoff preserves that horizontal support. Thus
\Cref{red:fixed-time} applies to the localized kernel error in
\eqref{red:scale-kernel-error}, and its \(L^p\) cost is uniform.
Only the original input and test carry horizontal restrictions;
no spatial support is subsequently imposed on their convolutions.

\subsection{Narrow vertical bins}

Split the vertical cutoff into its positive and negative frequency
parts. Reflection \(y\mapsto-y\) interchanges the two cases and
replaces \(u\) by \(-u\), preserving all assumptions. We work
with positive frequencies, so \(c_\eta\le w\eta\le C_\eta\).

Suppose first that \(L\ge L_0\) and \(m\) is as in
\eqref{red:small-time-parameters}. Choose real smooth cutoffs
\(Q_j\), \(1\le j\le J\le C m\), supported in intervals of
length \(C/W\) centered at \(\eta_j\asymp w^{-1}\), with
bounded overlap and
\[
 \sum_{j=1}^J Q_j(\eta)^2=1
       \quad\text{on }\supp p(w\eta).
\]
One obtains them from a fixed smooth partition at spacing \(W^{-1}\)
by dividing by the square root of its sum of squares on a slightly
enlarged band. Their normalized profiles and all fixed derivatives
are uniform. In particular, their convolution kernels have bounded
\(L^1\) norms, and
\begin{equation}\label{red:bin-mikhlin}
 A(Q_j)\le C m^4.
\end{equation}
Let \(\widetilde Q_j\) be a slightly enlarged cutoff equal to one
on the spectrum of \(Q_jP_w\). These cutoffs again have the stated
bin width and normalized smoothness bounds.

Set \(F_j^0=Q_jP_wf\) and \(G_j^0=Q_jg\). The identity on the
band and self-adjointness of the \(Q_j\) give
\begin{equation}\label{red:bin-diagonalization}
 \langle\mathcal A P_wf,g\rangle
   =\sum_j\langle\mathcal A\widetilde Q_jF_j^0,G_j^0\rangle
      +\sum_j\langle[\mathcal A,Q_j]F_j^0,g\rangle.
\end{equation}
Indeed, expand \(P_w=\sum_jQ_j^2P_w\), commute one \(Q_j\)
past \(\mathcal A\), and use
\(\widetilde Q_jF_j^0=F_j^0\). This identity does not require
\(\mathcal A\) to preserve individual bins.

On localized inputs and tests, Equations~\eqref{red:composition-bound},
\eqref{red:scale-first-moment}, and \eqref{red:bin-mikhlin} imply
\[
 \abs{\langle[\mathcal A,Q_j]F_j^0,g\rangle}
       \le C m^4\tau\norm{F_j^0}_2\norm g_2.
\]
In applying the composition estimate, only \(|t|\le b_0\) are
tested, as explained above. Bounded frequency overlap gives
\(\sum_j\norm{F_j^0}_2^2\le C\norm f_2^2\), so the total
error in \eqref{red:bin-diagonalization} is at most
\begin{equation}\label{red:bin-error}
 C m^5\tau\norm f_2\norm g_2.
\end{equation}
The harmless integer exponent five includes the factor
\(\sqrt J\). More generally, any fixed multiplier seminorm choice
would give \(C m^C\tau\), and for every fixed \(C,N\),
\begin{equation}\label{red:bin-error-small}
 m^C\tau
    \le C\exp\bigl(C' L^{0.9}-L^{0.97}\bigr)
    =O(L^{-N}).
\end{equation}
The large-\(L\) threshold in this statement is absolute after the
fixed exponents are chosen.

Let \(\mathcal H=\C^J\) with its Euclidean norm, and regard
\(F^0=(F_j^0)_j\), \(G^0=(G_j^0)_j\) as \(\mathcal H\)-valued
functions. For some fixed power \(Q=Cm^C\), their norms obey
\begin{equation}\label{red:bin-vector-norms}
 \begin{aligned}
 \norm{F^0}_{L^\infty(\mathcal H)}&\le Q,&
 \norm{F^0}_{L^1(\mathcal H)}&\le Q|F|,&
 \norm{F^0}_{L^2(\mathcal H)}^2&\le C|F|,\\
 \norm{G^0}_{L^\infty(\mathcal H)}&\le Q,&
 \norm{G^0}_{L^1(\mathcal H)}&\le Q|E|,&
 \norm{G^0}_{L^2(\mathcal H)}^2&\le C|E|.
 \end{aligned}
\end{equation}
For the first two columns, apply the uniform scalar kernel bounds
and sum over at most \(Cm\) components. For the last column, use
Plancherel and bounded overlap. No compact support in \(y\) is
asserted after the projections. Horizontal support in \(I\) is
preserved.

When \(m=1\), we do not perform this commutation or project the
test. Instead retain the original positive-frequency band cutoff
on the input and use one scalar component. It fits in a vertical
frequency interval centered at a fixed \(\eta_1\asymp w^{-1}\)
of length \(C/w=C/W\). Its positive annular support is retained.

\begin{lemma}[Amplitude levels after bin projection]
\label{red:bin-levels}
Let \(h:\R^2\to\mathcal H\) satisfy
\[
 \norm h_\infty\le Q,\qquad
 \norm h_1\le Q M_0,\qquad
 \norm h_2^2\le C_0M_0,
\]
where \(M_0>0\) and \(Q\le C\exp(C L^{0.9})\).
Write
\[
 A_i=\{2^{i-1}<\norm{h(\cdot)}_{\mathcal H}\le2^i\},
 \qquad a_i=2^i\sqrt{|A_i|/M_0}.
\]
For sufficiently large absolute \(L\),
\begin{equation}\label{red:level-weight-bounds}
 \sum_i a_i\le C\sqrt{1+L},\qquad
 \sum_{\{i:\,|\log(|A_i|/M_0)|>L/4\}}a_i
       \le C e^{-cL}.
\end{equation}
Empty level sets are omitted. The constants depend on the fixed
constants in the hypotheses, not on the dimension of \(\mathcal H\).
\end{lemma}

\begin{proof}
The amplitude upper bound implies \(2^i\le2Q\) on every nonempty
level. If \(|A_i|<e^{-L/4}M_0\), then
\(a_i\le2^i e^{-L/8}\), and summing the geometric series in
\(i\) gives \(C Qe^{-L/8}\).
If \(|A_i|>e^{L/4}M_0\), the \(L^1\) bound gives
\[
 2^i\le2Qe^{-L/4},\qquad
 a_i^2=4^i|A_i|/M_0\le2Q\,2^i.
\]
Summing the square-root geometric series over the allowed amplitudes
again gives \(C Qe^{-L/8}\). Since
\(\log Q=O(L^{0.9})\), both tails are \(O(e^{-cL})\).

The remaining level sets occupy at most \(C(1+L)\) dyadic measure
classes. Within a class, the geometric estimate in the proof of
\Cref{red:restricted-strong} bounds the sum of \(a_i\) by a
constant times its square sum. Cauchy--Schwarz over the classes and
\(\sum_i a_i^2\le4C_0\) give the first bound in
\eqref{red:level-weight-bounds}. The already bounded tails can be
included in that bound.
\end{proof}

Apply \Cref{red:bin-levels} to the two vectors in
\eqref{red:bin-vector-norms}. Decomposing them into amplitude levels
produces vector inputs and tests bounded by indicators of new sets.
The extra cutoffs \(\widetilde Q_j\) in
\eqref{red:bin-diagonalization} remain part of the operator after
this decomposition; thus no band support is being assumed for an
individual amplitude slice. If both new sets have retained measures,
then their logarithmic disparity is in
\begin{equation}\label{red:new-disparity}
 \left[\frac34L,\frac54L\right]
      \subset\left[\frac12L,2L\right].
\end{equation}
Indeed, the two changes of logarithmic measure are each at most
\(L/4\), and the logarithm of the square-root ratio uses one half
of their difference. The product of the two sums of normalized
amplitude weights is \(O(1+L)\).

\subsection{The admissible packet forms}

We next describe the packet model with all geometric and analytic
uniformities that will be used. Fix \(w\in2^{\Z}\), \(m\ge1\),
\(W=mw\), a finite dyadic scale set
\(\mathcal L\subset[w,\tau]\), and at most \(Cm\) bin centers
\(\eta_j\asymp w^{-1}\). The scalar case \(m=1\) just described
is included. Constants in this definition are fixed throughout.

\begin{definition}\label{red:packets}
A tile \(s\) has a length \(\ell_s\in\mathcal L\), a bounded
slope \(\theta_s\), a center \(z_s=(x_s,y_s)\), and precision
\(d_s=w/\ell_s\). Its box is
\[
 R_s=\left\{(x,y): |x-x_s|\le\frac{\ell_s}{2},\quad
       |(y-y_s)-\theta_s(x-x_s)|\le\frac W2\right\},
 \qquad |R_s|=\ell_s W.
\]
A \(t\)-box is defined by the same formula with slope \(t\).
The nominal horizontal length need not be the Euclidean long side
when \(W>\ell_s\).

At a fixed scale, slope labels have bounded multiplicity on every
interval of length \(d_s\). For each scale and label, the centers
lie in one translated box lattice
\[
 z_0+\{(n\ell_s,\theta_s n\ell_s+kW):n,k\in\Z\}.
\]
The lattice and center of a tile are common to every bin component.
An admissible collection is any finite subcollection of such lattices.

Its analyzing and synthesizing waves, denoted
\(\Phi_{s,j}\) and \(\Psi_{s,j}\), are bounded-normalized
modulated Schwartz functions adapted to \(R_s\). More precisely,
each can be written as
\[
 e^{i\eta_j((y-y_s)-\theta_s(x-x_s))}
 \phi_{s,j}\left(
       \frac{x-x_s}{\ell_s},
       \frac{(y-y_s)-\theta_s(x-x_s)}{W}\right),
\]
where, for all fixed nonnegative integers \(a,b,N\),
\begin{equation}\label{red:packet-seminorms}
 |\partial_X^a\partial_Y^b\phi_{s,j}(X,Y)|
       \le C_{a,b,N}(1+|X|+|Y|)^{-N}.
\end{equation}
The constants are uniform in \(s,j,w,m\). Both waves also have
Fourier support in
\begin{equation}\label{red:packet-spectra}
 |\eta-\eta_j|\le C W^{-1},\qquad
 c_\eta'\le w\eta\le C_\eta',\qquad
 |\xi+\theta_s\eta|\le c_*\ell_s^{-1},
\end{equation}
with fixed positive \(c_\eta',C_\eta'\). The constant \(c_*\)
is chosen sufficiently small compared with the selector separation
constants below and the fixed annular support of \(\psi\).

The scalar selector \(b_{s,j}\) is smooth, with
\begin{equation}\label{red:selector-assumptions}
 \supp b_{s,j}\subset
    \{t:c_b d_s\le|t-\theta_s|\le C_b d_s\},
 \qquad d_s^r\norm{b_{s,j}^{(r)}}_\infty\le C_r
\end{equation}
for every fixed \(r\ge0\), where \(0<c_b<C_b\) are absolute.
Finally, the collection uses one angular dyadic grid. Each tile has
an interval \(\Omega_s\) in that grid, of length
\(C_\Omega d_s\), where \(C_\Omega\) is a fixed large dyadic
number, containing both the wave slope support
\(\{-\xi/\eta:(\xi,\eta)\in\supp\widehat\Phi_{s,j}
\cup\supp\widehat\Psi_{s,j}\}\) and every selector support, with
a margin at least \(c_\Omega d_s\), where \(c_\Omega>0\) is fixed.
The interval is common to all components. Its level depends only on
\(\ell_s\).
\end{definition}

The associated form is
\begin{equation}\label{eq:src-3}
 \Lambda_{\mathcal S}(f,g)
   =\sum_{s\in\mathcal S}|R_s|^{-1}\sum_j
       \langle f_j,\Phi_{s,j}\rangle
       \langle b_{s,j}(u)\Psi_{s,j},g_j\rangle.
\end{equation}
It is linear in \(f\) and conjugate-linear in \(g\). The definition
makes sense for an arbitrary measurable bounded field \(u\). In
particular, the ordinary scalar and vector estimates in the next
section do not require spatial Lipschitz regularity of \(u\).
The stronger estimates later in the proof use the Lipschitz
hypothesis from \eqref{red:local-operator}.

\begin{proposition}[Packet representation]\label{red:packet-representation}
The diagonal form
\[
 \sum_j\langle\mathcal A_{u,\mathcal L}
                         \widetilde Q_j f_j,g_j\rangle
\]
is a sum, with absolutely summable scalar coefficients, of averages
of forms \eqref{eq:src-3}. Each form satisfies
\Cref{red:packets}, uniformly in \(w,m,\mathcal L\), and has one
common angular grid and common spatial lattices across components.
The same assertion holds for \(m=1\) with the original
positive-frequency band cutoff in place of \(\widetilde Q_1\).
The sums can be obtained as limits of finite packet collections.
\end{proposition}

\begin{proof}
We give the construction, including the normalization and the
separation of the different series forms. At each length \(\ell\),
partition the frequency slope \(-\xi/\eta\) smoothly into
intervals of width \(c\,d\), \(d=w/\ell\), with centers
\(\theta\) on a lattice of spacing a fixed small multiple of
\(d\). Take \(c\) so small that on every such frequency piece
\(\abs{\ell(\xi+\theta\eta)}\le c_*/4\).
Introduce normalized coordinates
\begin{equation}\label{red:normalized-symbol}
 X'=\ell(\xi+\theta\eta),\qquad
 Y'=W(\eta-\eta_j),\qquad
 U'=(u-\theta)/d.
\end{equation}
Then
\begin{equation}\label{red:symbol-identity}
 \ell(\xi+u\eta)=X'+(w\eta)U',\qquad
 w\eta=w\eta_j+Y'/m.
\end{equation}
The vertical cutoff restricts \(Y'\) to a fixed compact interval
and \(w\eta\) to a fixed positive annulus. On the support of
\(\psi\), its argument is bounded away from zero. In view of
\(\abs{X'}\le c_*/4\), this forces
\(c_b'\le|U'|\le C_b'\), with fixed positive bounds. Enlarge
these bounds slightly, retaining separation from zero. Since
\(d\le1\) and \(\abs u\le1\), only bounded values of
\(\theta\) can contribute.

In the coordinates \eqref{red:normalized-symbol}, the product of
the annular multiplier, the slope partition, and the vertical cutoff
is smooth on a fixed compact coordinate region, with all fixed
derivatives bounded uniformly. For example the slope partition is
a smooth function of
\(-X'/(w\eta)\), and \eqref{red:symbol-identity} shows that
its derivatives are uniform because \(w\eta\) stays away from
zero and \(m\ge1\). Choose slightly enlarged smooth plateaus in
\(X',Y',U'\), with the \(X'\)-plateau supported in
\(|X'|<c_*\), the \(U'\)-plateau supported away from zero, and
an additional vertical plateau keeping \(w\eta\) in a positive
annulus. Fourier series on a larger fixed coordinate box now give
\[
 \sum_{\nu\in\Z^3}
       a_{\nu,\ell,\theta,j}
       A_{\nu,\ell,\theta,j}(\xi,\eta)
       B_{\nu,\ell,\theta,j}((u-\theta)/d),
\]
where \(A_\nu\) is a localized exponential in \((X',Y')\),
\(B_\nu\) is a localized exponential in \(U'\), and for every
fixed \(N\),
\begin{equation}\label{red:fourier-series-decay}
 \sup_{\ell,\theta,j,w,m}|a_{\nu,\ell,\theta,j}|
          \le C_N(1+|\nu|)^{-N}.
\end{equation}
This follows by integration by parts on the fixed coordinate box.
Multiplication by the plateaus restores the original symbol exactly;
their supports give the strict spectral and selector restrictions
in \Cref{red:packets}.

Factor the frequency function \(A_\nu\) into two smooth frequency
functions by inserting a plateau equal to one on its support.
Both factors retain \(|X'|<c_*\), a vertical interval of length
\(C/W\), and a positive annular vertical support. Their inverse
Fourier transforms are modulated Schwartz kernels at the scales
\(\ell,W\). Their fixed seminorms, and those of the selector
factor, grow at most polynomially in \(|\nu|\). The coefficients
in \eqref{red:fourier-series-decay} absorb these losses uniformly.
One can do this simultaneously for every fixed seminorm as follows.
Let
\(A_\nu^*=\sup_{\ell,\theta,j,w,m}|a_{\nu,\ell,\theta,j}|\)
and omit zero values. Extract the common coefficient
\(c_\nu=(A_\nu^*)^{1/4}\), multiply each of the two frequency
factors by \((A_\nu^*)^{1/4}\), and multiply the selector factor
by \(a_{\nu,\ell,\theta,j}/(A_\nu^*)^{3/4}\).
The product is unchanged. Each of the three factors now has uniform
fixed seminorms, by \eqref{red:fourier-series-decay}, and
\(\sum_\nu c_\nu<\infty\). In particular, no derivative order
in this construction grows with \(L\) or a truncation parameter.

For clarity, consider one such factored term
\(b(u)A(D)B(D)\). Write \(K_A,K_B\) for its two convolution
kernels and \(|R|=\ell W\). At an intermediate convolution
center \(z\), put
\[
 \Psi_z(x)=|R|K_A(x-z),\qquad
 \Phi_z(y)=|R|\overline{K_B(z-y)}.
\]
Then the contribution to the bilinear form is exactly
\begin{equation}\label{red:continuous-packets}
 \frac1{|R|^2}\int_{\R^2}
       \langle f,\Phi_z\rangle
       \langle b(u)\Psi_z,g\rangle\dd z.
\end{equation}
The conjugation in \(\Phi_z\) keeps its Fourier support in the
same positive-frequency region as that of \(\Psi_z\).
Both waves have the carrier \(\eta_j(-\theta,1)\) and satisfy
\eqref{red:packet-seminorms}--\eqref{red:packet-spectra}.

Let \(\mathcal Q\) be a fundamental cell of the box lattice,
of area \(|R|\). For any integrable function \(F\),
\[
 \int_{\R^2}F(z)\dd z
    =\int_{\mathcal Q}\sum_{n,k\in\Z}
       F\bigl(z+(n\ell,\theta n\ell+kW)\bigr)\dd z.
\]
Applying this identity to \eqref{red:continuous-packets} and
normalizing the average over \(\mathcal Q\) leaves precisely
the factor \(|R|^{-1}\) in \eqref{eq:src-3}. Use the same
translation for every \(j\) at this scale and slope. There are
only finitely many lengths and relevant slope labels, so a finite
product of these normalized translation measures supplies a common
average for their complete sum.

We also need one angular grid per form. The union of the wave slope
supports and all selector supports for a label \(\theta\) is
contained in \([\theta-Cd,\theta+Cd]\), with \(C\) fixed and
independent of \(j\). Choose a sufficiently large fixed dyadic
multiple \(H d\). Among three adjacent dyadic grids, this interval
is contained in a cell of length \(H d\) with a margin comparable
to \(d\). One may see this directly from the three grids
\[
 \bigl\{2^{-k}([0,1)+n+(-1)^k\alpha):n,k\in\Z\bigr\},
       \qquad \alpha\in\{0,1/3,2/3\}:
\]
they are nested, and their boundary positions at any fixed scale
are spaced by one third of that scale. Taking \(H\) large gives
the claimed containment and margin. This is the usual adjacent-grid
principle; see \cite[Theorem~3.1 and Remark~3.2]{LernerNazarov2019}.
Assign each label to one suitable grid and split the sum into the
three resulting forms. Since \(d=w/\ell\) is dyadic, the chosen
cell level depends only on \(\ell\).

At a fixed length and label, the lattice sums converge absolutely
in a bilinear pairing of \(L^2\) functions. Indeed, envelope
decay and Cauchy--Schwarz give
\[
 \sum_z |R|^{-1}|\langle f,\Phi_z\rangle|^2
       \le C\norm f_2^2,
 \qquad
 \sum_z |R|^{-1}|\langle b_z(u)\Psi_z,g\rangle|^2
       \le C\norm g_2^2.
\]
For example, bound each squared pairing by \(C|R|\) times the
integral of \(|f|^2\) against a decaying box envelope, and then
sum those envelopes over the lattice. Their sum is uniformly
bounded. The same argument works with the bounded selector and
with vector components, by a further Cauchy--Schwarz inequality in
\(j\). Since the length and label sets are finite, the continuous
representation is justified by these bounds and the summability in
\(\nu\), first on regular inputs and then on \(L^2\).
Exhausting each lattice by finite subsets gives the claimed finite
packet approximations.

For each fixed series index and angular grid, all scales and labels
are kept in one form before its estimate is applied. Different series
indices remain separate forms and are summed only using
\(\sum_\nu c_\nu<\infty\). In particular, their copies of the
same spatial lattice are never combined into a collection with
unbounded multiplicity. This completes the representation.
\end{proof}

\subsection{The precise remaining model estimates}

The following proposition states the forward dependency used to
finish the band theorem. It also records exactly how much saving
survives the bin and amplitude reductions.

\begin{proposition}\label{red:model-suffices}
Suppose that the following bounds hold, uniformly for every admissible
finite form \(\Lambda_{\mathcal S}\) in
\Cref{red:packets}, with the fixed packet constants there.
The field \(u:\R^2\to\R\) is bounded and measurable, the sets
\(F,E\subset\R^2\) are measurable with positive finite measure,
and the input and test satisfy
\(\norm{f(z)}_{\mathcal H}\le\one_F(z)\),
\(\norm{g(z)}_{\mathcal H}\le\one_E(z)\).
\begin{enumerate}[label=\textup{(\roman*)}]
\item For \(m\ge1\), there is a fixed exponent \(C_1\) such that
\begin{equation}\label{red:model-basic}
 |\Lambda_{\mathcal S}(f,g)|
     \le C(1+\log(2m))^{C_1}\sqrt{|F||E|}.
\end{equation}
\item For \(L\ge L_0\), with \(\tau,m,W\) as in
\eqref{red:small-time-parameters}, assume in addition that
\(\abs u\le1\), \(\Lip(u)\le C_{\rm Lip}\), and
\(F,E\subset I\times\R\) for an interval \(I\) with
\(\abs I\le b_0\). Whenever
\(L/2\le\log\sqrt{|E|/|F|}\le2L\),
\begin{equation}\label{red:model-improved}
 |\Lambda_{\mathcal S}(f,g)|
       \le C L^{-20}\sqrt{|F||E|}.
\end{equation}
\item For \(m=1\), one scalar component, and \(|E|\le|F|\),
there is a fixed \(c>0\) such that
\begin{equation}\label{red:model-reverse}
 |\Lambda_{\mathcal S}(f,g)|
       \le C\sqrt{|F||E|}(|E|/|F|)^c.
\end{equation}
\end{enumerate}
Then \Cref{thm:band-restricted} holds. The constants remain uniform
in \(w\), \(\epsilon\), and the number of retained scales.
\end{proposition}

\begin{proof}
We apply the assumed bounds to the local field and sets of
\Cref{thm:band-restricted}, for which \(\abs u\le1\),
\(\Lip(u)\le C_{\rm Lip}\), and \(F,E\subset I\times\R\) with
\(\abs I\le b_0\).
Suppose first that \(L=\log\sqrt{|E|/|F|}\ge L_0\).
The time and scale reductions give exponentially small normalized
errors by \eqref{red:restricted-error}. The bin error is smaller
than any fixed power of \(L^{-1}\), by
\eqref{red:bin-error-small}. It remains to estimate the diagonal
form in \eqref{red:bin-diagonalization}.

Apply the amplitude decompositions from \Cref{red:bin-levels} to
\(F^0,G^0\). The projections defining these vectors preserve horizontal
support. After omitting levels of measure zero, their level sets have
positive finite measure by \eqref{red:bin-vector-norms} and lie in
\(I\times\R\); the normalized slices satisfy the stated indicator
bounds. First use finite decompositions. Each pair of slices
has a packet representation by \Cref{red:packet-representation};
its cutoff is \(\widetilde Q_j\), so the slice itself need not
have narrow spectrum. Summing the common series coefficients and
averaging the translations preserve each assumed model bound.
Pairs for which both measures are retained satisfy
\eqref{red:new-disparity}; apply \eqref{red:model-improved} to
them. Their total contribution, divided by
\(\sqrt{|F||E|}\), is at most
\[
 C L^{-20}(1+L)\le C L^{-19}.
\]
If either measure is discarded, use \eqref{red:model-basic}.
The sum of those normalized weights is at most
\(C\sqrt{1+L}\,e^{-cL}+Ce^{-2cL}\); the factor
\((1+\log(2m))^{C_1}\) is a fixed power of \(L\).
These terms are therefore exponentially small. The estimates are
uniform for the finite amplitude sums. Their tails converge
absolutely by the same bounds, so they apply to the full diagonal
form. All errors and the main term are bounded by
\(C(1+L)^{-3}\sqrt{|F||E|}\), as required.

If \(L\) ranges over a fixed bounded interval, use \(m=1\).
No bin projections or their amplitude losses are needed. The time
and scale errors have bounded operator norms, and
\eqref{red:model-basic}, with the scalar packet representation,
bounds the remaining form. The factor
\((1+|\log(|F|/|E|)|)^{-3}\) is bounded below by a positive
absolute constant on this interval.

Finally suppose that \(L\) is large and negative. Again use
\(m=1\), retain the original scalar input and test, and apply
\eqref{red:model-reverse}. It gives a positive power of
\(|E|/|F|\). The errors give the power \((|E|/|F|)^{1/4}\)
by the fixed choice \(p=4\) in
\eqref{red:restricted-error}. Every fixed positive power of this
ratio is bounded by a constant times
\((1+|\log(|F|/|E|)|)^{-3}\). This proves the remaining case.

All arguments include an arbitrary positive inner cutoff. If
\(\epsilon>w\), the lower retained scale is raised to a dyadic
scale comparable to \(\epsilon\), which simply selects a finite
subcollection of the allowed scales. If it exceeds \(\tau\),
the middle form is empty and the error estimates suffice. If
\(w\ge\tau\), the short-time and horizontal bounds already
cover the whole remaining integral. None of the bounds used here
depends on how many scales lie between the endpoints.
\end{proof}

The basic and reverse-disparity assertions of
\Cref{red:model-suffices} will follow from the ordinary tile
estimates, in particular \Cref{tile:prop-model}. The remaining
large-positive-disparity assertion is the improvement proved by the
popular-forest argument and assembled in
\Cref{assembly:all-model-estimates}. Thus the conditional use of
\Cref{thm:band-restricted} in \Cref{red:band-implies-main} is
separate from, and subsequent to, the proof of its model hypotheses.
